\documentclass{article}
\usepackage[utf8]{inputenc} % Para admitir tildes correctamente
\usepackage[spanish]{babel}

\usepackage{amsmath}
\usepackage{amssymb}
\usepackage{amsthm}
%\usepackage{thmtools}

\theoremstyle{plain}
\newtheorem{teorema}{Teorema}[section]
\newtheorem{lema}[teorema]{Lema}

\theoremstyle{definition}
\newtheorem{definicion}{Definición}[section]
\newtheorem{formula}[teorema]{Fórmula}

\begin{document}

\title{Mis Apuntes de Números Combinatorios}
\author{David Bernal}
\maketitle

\section{Introducción}
Los números combinatorios $\binom{n}{k}$ aparecen al estudiar los coeficientes que van apareciendo en:
\begin{align*}
(a+b)^0 &= 1 \\
(a+b)^1 &= a + b \\
(a+b)^2 &= a^2 + 2ab + b^2 \\
(a+b)^3 &= a^3 + 3a^2b + 3ab^2 + b^3 \\
(a+b)^4 &= a^4 + 4a^3b + 6a^2b^2 + 4ab^3 + b^4 \\
&\ \vdots
\end{align*}

Este desarrollo puede hacerse en cualquier anillo conmutativo.
Notemos que todos los sumandos de $(a+b)^n$ serán de la forma $a^{n-m}b^m$, y llamando $\binom{n}{m}$
al número que resulta de juntar todos los elementos iguales a $a^{n-m}b^m$, tenemos lo que se llama fórmula
del binomio de Newton:

\begin{equation}
(a+b)^n = \displaystyle\sum_{m=0}^n \binom{n}{m}a^{n-m}b^m.
\end{equation}

A los números $\binom{n}{m}$ se los llama números combinatorios, con $ 0 \le m \le n$. Y se calculan así:
\begin{teorema}
    Se verifica la siguiente igualdad para $ 0 \le m \le n$
    \[
    \binom{n}{m} = \frac{n!}{m!(n-m)!}
    \]
\end{teorema}

Antes de ver la demostración, veamos unas consecuencias sencillas de ésta identidad.
\begin{enumerate}
    \item Se tiene
    \[
    \binom{n}{m} = \binom{n}{n-m}.
    \]
    En efecto, $\binom{n}{m} = \frac{n!}{m!(n-m)!} = \frac{n!}{(n-(n-m))!(n-m)!} = \binom{n}{n-m}.$
    
    \item Tambien tenemos la siguiente igualdad para $0 < m < n$
    \[
    \binom{n+1}{m+1} = \binom{n}{m} + \binom{n}{m+1}.
    \]
    Tenemos que $\binom{n+1}{m+1} = \frac{(n+1)!}{(m+1)!(n-m)!} = \frac{(m+1)+(n-m)}{m+1}\cdot\binom{n}{m} = \binom{n}{m} + \binom{n}{m+1}$
\end{enumerate}

\end{document}
